%% kapa: exam=1 round=35 year=2024 period=1 form=A subject=civil_law subject_ko=민법 seq=15 number=15
%% source: Q-Net 감정평가사 기출문제 (cst00309) artlSeq=5213229 file=raw/1008/5213229/1교시 문제지 원본.pdf
%% conversion: kordoc 4.24.51 -> tools/gen1.py (auto); review=proofread-sonnet; flags=none
\begin{question}{15}
甲은 토지거래허가구역 내에 있는 자신의 X토지에 대해 허가를 받을 것을 전제로 乙에게 매도하는 계약을 체결하였으나 아직 허가는 받지 않은 상태이다. 이에 관한 설명으로 옳지 않은 것은? (다툼이 있으면 판례에 따름)
\begin{choices}
  \item 乙은 甲에게 계약의 이행을 청구할 수 없다.
  \item 甲이 토지거래허가신청절차에 협력하지 않는 경우, 乙은 이를 이유로 계약을 해제할 수 있다.
  \item 토지거래허가구역 지정이 해제된 경우, 특별한 사정이 없는 한 위 매매계약은 확정적으로 유효하다.
  \item 甲과 乙이 토지거래허가를 받으면 위 매매계약은 소급해서 유효로 되므로 허가 후에 새로 매매계약을 체결할 필요는 없다.
  \item 甲의 사기에 의하여 위 매매계약이 체결된 경우, 乙은 토지거래허가를 신청하기 전이라도 甲의 사기를 이유로 매매계약을 취소할 수 있다.
\end{choices}
\end{question}
\answer{2}
